\documentclass[12pt]{article}
\usepackage[utf8]{inputenc}
\usepackage[T1]{fontenc}
\usepackage{hyperref}
\title{Umbindhamu Language Resources for K-12 Curriculum Development}
\author{Compiled by Claude Sonnet 5 (claude-sonnet-5) for Andr\'as Kornai}
\date{}
\begin{document}
\maketitle

\section{Introduction}
Umbindhamu is an alternative spelling of Umpithamu (ISO 639-3 umd), an Aboriginal language of Cape York Peninsula, Queensland. Sources report very few or no remaining fluent speakers. Search results mention two fluent speakers in the early 2000s, one of them Florrie Bassani. Classification is disputed.

\section{Primary Educational Resources}
The search found no primers, textbooks or school programmes.

\section{Linguistic and Reference Materials}
\begin{itemize}
\item Wikipedia article ``Umpithamu language'', which says Jean-Christophe Verstraete published \emph{A Dictionary of Umpithamu} in 2020: \url{https://en.wikipedia.org/wiki/Umpithamu_language}
\item Glottolog reference record (Laycock 1969): \url{https://glottolog.org/resource/reference/id/303408}
\end{itemize}

\section{Religious and Community Resources}
None found in the search.

\section{Digital and Audio Resources}
The State Library of Queensland language blog covers the related Umbuygamu language, not umd itself: \url{https://www.slq.qld.gov.au/blog/language-week-week-thirty-umbuygamu}

\section{Curriculum Development Context}
Very little is available. A dictionary exists, but I did not see it directly. Any revitalization curriculum would need to work with the descendant community and with the dictionary authors. No such contacts were found in the search.
\end{document}